\section*{ЗАКЛЮЧЕНИЕ}
\addcontentsline{toc}{section}{ЗАКЛЮЧЕНИЕ}

В результате выполнения выпускной квалификационной работы разработана веб-ориентированная программно-информационная система для управления проектами. Система рассчитана на небольшие команды, ведущие несколько параллельных программных проектов, и доступна пользователям через современный браузер без установки клиентского программного обеспечения.

В ходе работы проведен анализ предметной области управления проектами и сопоставлены функциональные возможности восьми существующих ИС~УП, доступных российским пользователям. Анализ показал, что ни одна из рассмотренных систем не объединяет в одной точке нативную диаграмму Ганта, поддержку четырех типов зависимостей сетевого планирования, русскоязычный интерфейс, бесплатный тариф для небольшой команды и устойчивый доступ из РФ. По результатам сравнения сформулированы требования к разрабатываемой системе.

Разработанная система реализует следующие основные возможности:
\begin{itemize}
    \item регистрацию и аутентификацию пользователей с сеансовым механизмом доступа;
    \item ведение реестра проектов и декомпозицию работ на задачи;
    \item построение диаграммы Ганта как основного представления календарного плана;
    \item интерактивное изменение дат задач и их порядка непосредственно на диаграмме;
    \item установление зависимостей между задачами с поддержкой четырех типов сетевого планирования: финиш-старт, старт-старт, финиш-финиш, старт-финиш;
    \item автоматическую проверку циклов в графе зависимостей и согласованности дат с типом связи;
    \item разграничение прав участников проекта по трем ролям: владелец, редактор, наблюдатель;
    \item управление составом участников проекта с поиском пользователей по логину и полному имени;
    \item настройку цветовой темы интерфейса в четырех режимах: системный, светлый, темный, пользовательский;
    \item адаптивный интерфейс, корректно отображающийся на устройствах с различными размерами экрана.
\end{itemize}

Для реализации системы выбрана трехуровневая клиент-серверная архитектура. Серверная часть реализована на платформе Django с библиотекой Django REST Framework, клиентская -- на библиотеке React. В качестве постоянного хранилища использована реляционная СУБД MySQL. Доступ к базе организован через ORM Django с параметризованными запросами, что исключает основную поверхность атак через внедрение SQL-кода. Аутентификация выполнена на стандартном механизме сеансов Django, изменяющие операции защищены CSRF-токеном. Пароли хранятся в виде хеша, открытый пароль нигде не сохраняется.

Системное тестирование подтвердило корректность работы реализованной функциональности по всем сценариям технического задания.

Основные результаты выпускной квалификационной работы:
\begin{enumerate}
    \item Проведен анализ предметной области управления проектами и сравнительный анализ восьми существующих ИС~УП, доступных российским пользователям. Обоснована потребность в собственной разработке.
    \item Сформировано техническое задание на разработку, описаны прецеденты использования системы по четырем ролям пользователей.
    \item Спроектирована архитектура программной системы, построена ER-модель данных, описана структура базы данных, разработан программный интерфейс между клиентской и серверной частями.
    \item Реализована серверная часть программной системы на Django и Django REST Framework с использованием СУБД MySQL.
    \item Реализована клиентская часть программной системы на React с интерактивной диаграммой Ганта, поддерживающей четыре типа зависимостей задач, перетаскивание полос и изменение порядка строк.
    \item Реализована ролевая модель разграничения прав участников проекта и управление составом команды.
    \item Реализована настройка цветовой темы интерфейса с сохранением палитры в учетной записи пользователя.
    \item Проведено системное тестирование, подтвердившее корректность работы реализованной функциональности.
\end{enumerate}

Все задачи, поставленные в начале выпускной квалификационной работы, успешно решены. Цель работы -- разработка программно-информационной системы для управления проектами -- достигнута.
